%% kapa: exam=2 round=34 year=2023 period=3 subject=compensation_law subject_ko=감정평가 및 보상법규 seq=2 number=2 points=30
%% source: Q-Net 감정평가사 기출문제 (cst00309) artlSeq=5212168 file=raw/1008/5212168/3교시 23년도 감정평가사 2차시험_감정평가 및 보상법규.pdf
%% conversion: kordoc 4.24.51 -> tools/gen2.py (auto); review=proofread-gpt-sol; flags=none
\begin{longquestion}{【문제 2】}
지적공부상 지목이 전(田)인 甲소유의 토지(‘이 사건 토지’라 함)는 면적이 2,000 m²이고, 이 중 330 m² 토지에 주택이 건축되어 있고 나머지 부분은 밭으로 사용되고 있다. 그럼에도 불구하고 A도 B시의 시장(이하 ‘B시장’이라 함)은 지목이 대(垈)인 1개의 표준지의 공시지가를 기준으로 토지가격비준표를 사용하여 2022. 5. 31. 이 사건 토지에 대하여 개별공시지가를 결정ㆍ공시하였다. B시장은 이 사건 토지에 대한 개별공시지가와 이의신청 절차를 甲에게 통지하였다. 다음 물음에 답하시오. (단, 각 물음은 상호독립적임) (30점) \par
\medskip\noindent 물음 1) 甲이 B시장의 개별공시지가결정이 위법ㆍ부당하다는 이유로 부동산 가격 공시에 관한 법령에 따른 이의신청을 거치지 않고 ｢행정심판법｣에 따른 취소심판을 제기할 수 있는지 여부와 이 사건 토지에 대한 개별공시지가결정의 위법성에 관하여 설명하시오. (15점) \par
\medskip\noindent 물음 2) 甲은 개별공시지가결정에 대하여 부동산 가격공시에 관한 법령에 따른 이의신청이나 ｢행정심판법｣에 따른 행정심판과 ｢행정소송법｣에 따른 행정소송을 제기하지 않았다. 그 후 B시장은 2022. 9. 15. 이 사건 토지에 대한 개별공시지가를 시가표준액으로 하여 재산세를 부과ㆍ처분하였다. 이에 甲은 2022. 12. 5. 이 사건 토지에 대한 개별공시지가결정의 하자를 이유로 재산세부과처분에 대하여 취소소송을 제기하였다. 甲의 청구가 인용될 수 있는지 여부에 관하여 설명하시오. (15점) \par
\begin{framedtext}
참조 조문 \par
<부동산 가격공시에 관한 법률> \par
제10조(개별공시지가의 결정ㆍ공시 등) ① 시장ㆍ군수 또는 구청장은 국세ㆍ지방세 등 각종 세금의 부과, 그 밖의 다른 법령에서 정하는 목적을 위한 지가산정에 사용되도록 하기 위하여 제25조에 따른 시ㆍ군ㆍ구부동산가격공시위원회의 심의를 거쳐 매년 공시지가의 공시기준일 현재 관할 구역 안의 개별토지의 단위면적당 가격(이하 “개별공시지가”라 한다)을 결정ㆍ공시하고, 이를 관계 행정기관 등에 제공하여야 한다. \par
② 제1항에도 불구하고 표준지로 선정된 토지, 조세 또는 부담금 등의 부과대상이 아닌 토지, 그 밖에 대통령령으로 정하는 토지에 대하여는 개별공시지가를 결정ㆍ공시하지 아니할 수 있다. 이 경우 표준지로 선정된 토지에 대하여는 해당 토지의 표준지공시지가를 개별공시지가로 본다. \par
<생략> \par
④ 시장ㆍ군수 또는 구청장이 개별공시지가를 결정ㆍ공시하는 경우에는 해당 토지와 유사한 이용가치를 지닌다고 인정되는 하나 또는 둘 이상의 표준지의 공시지가를 기준으로 토지가격비준표를 사용하여 지가를 산정하되, 해당 토지의 가격과 표준지공시지가가 균형을 유지하도록 하여야 한다. \par
⑤ 시장ㆍ군수 또는 구청장은 개별공시지가를 결정ㆍ공시하기 위하여 개별토지의 가격을 산정할 때에는 그 타당성에 대하여 감정평가법인등의 검증을 받고 토지소유자, 그 밖의 이해관계인의 의견을 들어야 한다. 다만, 시장ㆍ군수 또는 구청장은 감정평가법인등의 검증이 필요 없다고 인정되는 때에는 지가의 변동상황 등 대통령령으로 정하는 사항을 고려하여 감정평가법인등의 검증을 생략할 수 있다. \par
제11조(개별공시지가에 대한 이의신청) ① 개별공시지가에 이의가 있는 자는 그 결정ㆍ공시일부터 30일 이내에 서면으로 시장ㆍ군수 또는 구청장에게 이의를 신청할 수 있다. \par
② 시장ㆍ군수 또는 구청장은 제1항에 따라 이의신청 기간이 만료된 날부터 30일 이내에 이의신청을 심사하여 그 결과를 신청인에게 서면으로 통지하여야 한다. 이 경우 시장ㆍ군수 또는 구청장은 이의신청의 내용이 타당하다고 인정될 때에는 제10조에 따라 해당 개별공시지가를 조정하여 다시 결정ㆍ공시하여야 한다. \par
③ 제1항 및 제2항에서 규정한 것 외에 이의신청 및 처리절차 등에 필요한 사항은 대통령령으로 정한다. \par
<부동산 가격공시에 관한 법률 시행령> \par
제21조(개별공시지가의 결정 및 공시) ① 시장ㆍ군수 또는 구청장은 매년 5월 31일까지 개별공시지가를 결정ㆍ공시하여야 한다. 다만, 제16조제2항제1호의 경우에는 그 해 10월 31일까지, 같은 항 제2호의 경우에는 다음 해 5월 31일까지 결정ㆍ공시하여야 한다. \par
② 시장ㆍ군수 또는 구청장은 제1항에 따라 개별공시지가를 공시할 때에는 다음 각 호의 사항을 해당 시ㆍ군 또는 구의 게시판 또는 인터넷 홈페이지에 게시하여야 한다. \par
\textbf{1. 조사기준일, 공시필지의 수 및 개별공시지가의 열람방법 등 개별공시지가의 결정에 관한 사항} \par
\textbf{2. 이의신청의 기간ㆍ절차 및 방법} \par
③ 개별공시지가 및 이의신청기간 등의 통지에 관하여는 제4조제2항 및 제3항을 준용한다. \par
제22조(개별공시지가에 대한 이의신청) ① 법 제11조제1항에 따라 개별공시지가에 대하여 이의신청을 하려는 자는 이의신청서에 이의신청 사유를 증명하는 서류를 첨부하여 해당 시장ㆍ군수 또는 구청장에게 제출하여야 한다. \par
② 시장ㆍ군수 또는 구청장은 제1항에 따라 제출된 이의신청을 심사하기 위하여 필요할 때에는 감정평가법인등에게 검증을 의뢰할 수 있다. \par
\end{framedtext}
\end{longquestion}
